\chapter{引言}\label{chap:introduction}

随着机器人技术与智能控制方法的发展，机械臂在动态分拣、移动目标拦截、自动装配和非合作目标捕获等场景中的应用需求不断增加。这类任务和传统静态抓取相比存在明显差异：目标的位置、姿态和速度在执行过程中持续变化，机器人需要同时完成运动预测、末端对齐、夹爪闭合和后续运输。若观测来自视觉系统，还会受到测量噪声、延迟和丢帧等因素影响，使控制系统面临更高的不确定性。

传统基于模型的控制方法依赖明确的运动学模型和状态估计，具有结构清晰、可解释性强、约束处理直接等优点。在仿真真值观测下，模型预测控制器能够根据目标当前状态构造参考轨迹并完成稳定执行。然而在观测含噪的情况下，确定性控制器可能会把目标位姿的高频抖动直接传递到末端轨迹或关节命令中，导致抓取时机偏移、求解不稳定或动作震荡。与此相对，模仿学习通过专家示教数据直接学习从观测到动作的映射，有机会利用历史观测和动作块输出缓解部分噪声影响，但其性能高度依赖示教数据质量。

模仿学习在机器人操作任务中已经形成了行为克隆、分层模仿学习和生成式策略等多种路线。行为克隆把专家示教转化为监督学习问题，训练过程相对稳定，但在复杂动态场景中容易受分布偏移影响。分层模仿学习将高层计划和低层控制分开，降低长时程任务的学习难度。近年来扩散模型在生成建模中表现突出，被用于机器人动作块生成后，能够表达多峰动作分布并生成较平滑的动作序列\citep{ho2020denoising,song2020denoising,chi2023diffusion}。这些特点使扩散策略适合用于动态目标抓取中的短时参考轨迹生成。

本文围绕仿真中的动态立方体抓取任务展开研究，采用“算法专家生成示教，学习策略蒸馏专家行为”的技术路线。算法专家在仿真真值观测下通过有限状态机和模型预测控制生成示教轨迹，随后将成功回合编译为目标中心化的训练窗口，训练条件扩散策略根据历史相对观测和短时未来目标条件输出末端执行器动作块。运行时，扩散策略负责抓取阶段参考生成，MPC控制器负责将任务空间参考转化为可执行动作。

本文主要完成了以下工作。首先，构建了动态目标抓取仿真、数据采集、模型训练和批量评估工作流，支持真值观测和多类噪声观测配置。其次，采集并整理了中等难度任务的合成专家数据，原始数据包含12500个回合和2934103帧，编译后得到11909个成功回合对应的扩散策略训练窗口。再次，训练目标中心化的条件扩散策略，并在真值观测和含噪观测矩阵上与MPC专家进行比较。最后，通过200回合评估结果分析两类策略在不同运动分量和噪声类型下的优势与失效模式。

本文的组织结构如下：第二章介绍模仿学习、扩散策略和模型预测控制等相关工作；第三章说明仿真环境、专家数据、扩散策略和评估矩阵的具体设计；第四章展示实验结果并分析真值观测和含噪观测下的性能差异；第五章总结全文并讨论后续改进方向。
